\subsection{Why access, use and control can come apart}
\label{sec:framework}
Whether a household keeps a remittance in digital form depends on what it can do with those funds. A digital balance is useful if the recipient can access the account, pay the people or merchants she needs to pay, and resolve a problem when something goes wrong. Cash can be a cheaper or more reliable choice for particular obligations. So can assistance from a trusted person. These choices are not an ordered ladder from exclusion to inclusion. They depend on where the money is going and who can move it.

\paragraph{Usable access.} A payment rail can lower the cost of moving funds between participating accounts without changing whether a particular woman can open or use an account independently. A phone and proof of address in her own name may facilitate verification. Their absence need not prevent transactions arranged through someone else's account, and their presence does not show that a provider will accept the documentation. Under this mechanism, own-name conditions and reported use can move together, but the link also depends on help, provider practices and how the account is used. The distinction among resources, agency and achievements \citep{kabeer1999} explains why neither account access nor a completed transfer is a sufficient measure of control.

\paragraph{Relative costs and cash liquidity.} Compare the full cost of completing a payment digitally with the cost of taking funds out and paying another way. Quoted fees are only one part of that comparison. Exchange-rate spreads, trips to an agent, delays, failed attempts and verification can affect which route is preferable. A fast domestic transfer can be the cheapest first leg of an international transaction and still leave the sender relying on an intermediary at the border. Opportunities for risk sharing are a possible consequence of lower transaction costs. Evidence on mobile money in Kenya illustrates the difference between measuring adoption and observing changes in consumption after shocks \citep{jackSuri2014}. We observe reported uses and experiences, not the corresponding change in welfare.

\paragraph{Control and recourse.} Help from relatives, agents or acquaintances can make an otherwise unavailable transaction possible. Its implications differ if the user can choose whom to ask, keep credentials private, verify delivery and obtain help after a failure. Reports of precautions are not evidence that a remedy exists; an institutional complaint channel is not evidence of restitution. In the survey these are separate questions. The authored accounts help identify which distinctions matter in practice and cases where the same answer could reflect different arrangements.

None of these mechanisms implies that the two survey outcomes must be highly correlated or that a particular group of respondents should form a latent class. The survey measures reported conditions and uses. The interviews illuminate steps and constraints that a survey answer may combine. Policy and provider evidence identifies some of the rules through which those constraints can arise. Their different units and observation processes matter for what each can support \citep{adcockCollier2001,humphreysJacobs2023}.
